% ============================================================
%  GUÍA 7: APLICACIONES DE LA DERIVADA
%  Curso de Introducción a la Universidad — Precálculo y Cálculo I
%  Autor: Gabriel Astudillo
%  Sesiones cubiertas: 21, 22 y 23
% ============================================================

\documentclass[12pt, a4paper]{article}

% ── Paquetes ─────────────────────────────────────────────────

\usepackage{mathwizards-guia}
\mwguia{Guía 7 — Aplicaciones de la Derivada}{Gabriel Astudillo $\cdot$ Curso Preuniversitario}

\begin{document}
% ============================================================

% ── Portada / Título ─────────────────────────────────────────
\mwportada{Guía de Estudio 7}{Aplicaciones de la Derivada}{Máximos y Mínimos, Optimización, Razones de Cambio, Teorema de Rolle y TVM, Análisis de Gráficas}

\vspace{4pt}
\begin{cuadroinfo}
  \small
  \textbf{Presentación:} Esta guía muestra el poder de la derivada en situaciones reales. Aprenderás a encontrar valores extremos de funciones, resolver problemas de optimización, calcular razones de cambio relacionadas y analizar la forma de una gráfica. Cada sección incluye \textbf{ejercicios resueltos} y \textbf{propuestos}.
  \\[4pt]
  \textbf{Nivel de Dificultad:}\quad
  \facil{} Básico / Directo \quad
  \medio{} Intermedio \quad
  \dificil{} Desafiante
\end{cuadroinfo}

\vspace{8pt}

% ============================================================
\section{Máximos y Mínimos Locales}
% ============================================================

\subsection{Conceptos fundamentales}

\begin{cuadroteoria}
  \textbf{Valor máximo local:} $f(c)$ es máximo local si $f(c) \geq f(x)$ para todo $x$ cercano a $c$.

  \textbf{Valor mínimo local:} $f(c)$ es mínimo local si $f(c) \leq f(x)$ para todo $x$ cercano a $c$.

  \textbf{Punto crítico:} Un punto $c$ del dominio de $f$ donde $f'(c) = 0$ o $f'(c)$ no existe.

  \textbf{Teorema (condición necesaria):} Si $f$ tiene un extremo local en $c$ y $f'(c)$ existe, entonces $f'(c) = 0$. Los extremos locales ocurren en puntos críticos.
\end{cuadroteoria}

\begin{cuadrosolucion}
  \textbf{Prueba de la primera derivada:}
  \begin{itemize}[leftmargin=*]
    \item Si $f'$ cambia de $+$ a $-$ al pasar por $c$: máximo local en $c$.
    \item Si $f'$ cambia de $-$ a $+$ al pasar por $c$: mínimo local en $c$.
    \item Si $f'$ no cambia de signo: ni máximo ni mínimo.
  \end{itemize}
\end{cuadrosolucion}

\begin{cuadroextra}
  \textbf{Prueba de la segunda derivada:}
  Sea $c$ un punto crítico donde $f'(c) = 0$:
  \begin{itemize}[leftmargin=*]
    \item Si $f''(c) < 0$: $f$ es cóncava hacia abajo en $c \;\Rightarrow\; f(c)$ es un \textbf{máximo local}.
    \item Si $f''(c) > 0$: $f$ es cóncava hacia arriba en $c \;\Rightarrow\; f(c)$ es un \textbf{mínimo local}.
    \item Si $f''(c) = 0$: la prueba \textbf{no es concluyente} (usar la primera derivada).
  \end{itemize}
\end{cuadroextra}

\newpage

\begin{cuadroadvertencia}
  \textbf{Teorema de Fermat:} Si $f$ es continua en $[a, b]$, entonces $f$ alcanza un valor máximo y un valor mínimo absolutos en ese intervalo. Para encontrarlos:
  \begin{enumerate}
    \item Halla los puntos críticos en $(a, b)$.
    \item Evalúa $f$ en los puntos críticos y en los extremos $a$ y $b$.
    \item El mayor valor es el máximo absoluto; el menor, el mínimo absoluto.
  \end{enumerate}
\end{cuadroadvertencia}

\subsection{Ejercicios resueltos}

\textbf{Ejemplo R1.} Encuentra los puntos críticos de $f(x) = x^3 - 3x^2 + 2$.

\begin{cuadrosolucion}
  \textbf{Solución:}
  \begin{enumerate}
    \item $f'(x) = 3x^2 - 6x = 3x(x - 2)$.
    \item $f'(x) = 0 \Rightarrow 3x(x - 2) = 0 \Rightarrow x = 0$ o $x = 2$.
  \end{enumerate}
  Puntos críticos: $x = 0$ y $x = 2$.
\end{cuadrosolucion}

\textbf{Ejemplo R2.} Clasifica los puntos críticos de $f(x) = x^3 - 3x^2$ usando la prueba de la primera derivada.

\begin{cuadrosolucion}
  \textbf{Solución:} $f'(x) = 3x(x - 2)$. Los puntos críticos son $x = 0$ y $x = 2$. Analizamos signos:
  \begin{center}
    \small
    \begin{tabular}{lccc}
      \toprule
      \textbf{Intervalo} & $(-\infty, 0)$ & $(0, 2)$    & $(2, \infty)$ \\
      \midrule
      Signo de $f'$      & $+$            & $-$         & $+$           \\
      Comportamiento     & creciente      & decreciente & creciente     \\
      \bottomrule
    \end{tabular}
  \end{center}
  En $x = 0$: $f'$ cambia de $+$ a $-$, hay máximo local $f(0) = 0$.

  En $x = 2$: $f'$ cambia de $-$ a $+$, hay mínimo local $f(2) = -4$.
\end{cuadrosolucion}

\textbf{Ejemplo R3.} Encuentra los extremos absolutos de $f(x) = x^2 - 4x + 5$ en $[0, 5]$.

\begin{cuadrosolucion}
  \textbf{Solución:}
  \begin{enumerate}
    \item Derivada: $f'(x) = 2x - 4$. Cero: $x = 2$ (en el intervalo).
    \item Evaluamos:
          \begin{itemize}
            \item $f(0) = 5$
            \item $f(2) = 4 - 8 + 5 = 1$
            \item $f(5) = 25 - 20 + 5 = 10$
          \end{itemize}
    \item Mínimo absoluto: $1$ (en $x = 2$). Máximo absoluto: $10$ (en $x = 5$).
  \end{enumerate}
\end{cuadrosolucion}

\newpage

\textbf{Ejemplo R4.} Encuentra los puntos críticos de $f(x) = \sqrt[3]{(x - 1)^2}$.

\begin{cuadrosolucion}
  \textbf{Solución:} Reescribimos $f(x) = |x - 1|^{2/3}$. La derivada:
  $$f'(x) = \frac{2}{3}(x - 1)^{-1/3} = \frac{2}{3\sqrt[3]{x - 1}}.$$
  La derivada nunca es cero, pero no existe en $x = 1$. Por tanto, $x = 1$ es punto crítico (donde la función tiene un pico).
\end{cuadrosolucion}

\subsection{Ejercicios propuestos}

\begin{enumerate}[leftmargin=*]
  \item \facil{} Encuentra los puntos críticos de $f(x) = x^2 - 6x + 10$.

  \item \facil{} Encuentra los puntos críticos de $f(x) = x^4 - 4x^2$.

  \item \facil{} Encuentra los extremos absolutos de $f(x) = x^2$ en $[-2, 1]$.

  \item \facil{} Encuentra los extremos absolutos de $f(x) = x + 1$ en $[0, 4]$.

  \item \medio{} Clasifica los puntos críticos de $f(x) = 2x^3 - 3x^2 - 36x$.

  \item \medio{} Encuentra los extremos absolutos de $f(x) = x^3 - 3x + 1$ en $[0, 3]$.

  \item \medio{} Encuentra los puntos críticos de $f(x) = \dfrac{x}{x^2 + 1}$.

  \item \medio{} Determina los máximos y mínimos locales de $f(x) = \sin x$ en $[0, 2\pi]$.

  \item \dificil{} Encuentra los valores de $a$ para que $f(x) = ax^3 - 6x^2$ tenga un punto crítico en $x = 2$.

  \item \dificil{} Determina los extremos absolutos de $f(x) = x\sqrt{9 - x^2}$ en $[-3, 3]$.

  \item \dificil{} Encuentra los extremos absolutos de $f(x) = \dfrac{x - 1}{x + 1}$ en $[0, 4]$.

  \item \dificil{} Prueba que $f(x) = x^3 + x + 1$ tiene exactamente un punto crítico (y que es un mínimo local). ¿Cuál es?
\end{enumerate}

\newpage

% ============================================================
\section{Problemas de Optimización}
% ============================================================

\subsection{Metodología para resolver problemas de optimización}

\begin{cuadroteoria}
  \textbf{Pasos recomendados:}
  \begin{enumerate}
    \item \textbf{Identifica} la cantidad a maximizar o minimizar (llámala $Q$).
    \item \textbf{Expresa} $Q$ en términos de una sola variable, usando las restricciones del problema (puentes geométricos como perímetros, áreas, volúmenes).
    \item \textbf{Determina} el dominio realista de la variable.
    \item \textbf{Deriva} $Q$ respecto a la variable.
    \item \textbf{Iguala} la derivada a cero y resuelve.
    \item \textbf{Verifica} si se trata de un máximo o mínimo (primera o segunda derivada), y comprueba los extremos del dominio.
  \end{enumerate}
\end{cuadroteoria}

\begin{cuadroadvertencia}
  \textbf{Recuerda:} en problemas de la vida real, los valores negativos o fuera del rango no tienen sentido. Verifica siempre la plausibilidad de la solución.
\end{cuadroadvertencia}

\subsection{Ejercicios resueltos}

\textbf{Ejemplo R1.} Un granjero quiere cercar un campo rectangular de área $1200$ m$^2$ usando la menor cantidad de cerca posible. ¿Cuáles deben ser las dimensiones?

\begin{cuadrosolucion}
  \textbf{Solución:} Sea $x$ el largo e $y$ el ancho. Condición: $xy = 1200 \Rightarrow y = \frac{1200}{x}$. Perímetro: $P = 2x + 2y = 2x + \frac{2400}{x}$, con $x > 0$.
  \begin{itemize}
    \item Derivamos: $P'(x) = 2 - \dfrac{2400}{x^2}$.
    \item $P'(x) = 0 \Rightarrow 2 = \dfrac{2400}{x^2} \Rightarrow x^2 = 1200 \Rightarrow x = \sqrt{1200} = 20\sqrt{3} \approx 34{,}64$ m.
    \item $P''(x) = \dfrac{4800}{x^3} > 0$ siempre, así que es mínimo.
    \item $y = \dfrac{1200}{20\sqrt{3}} = \dfrac{60}{\sqrt{3}} = 20\sqrt{3}$ m.
  \end{itemize}
  Las dimensiones óptimas son $20\sqrt{3} \times 20\sqrt{3}$ m (un cuadrado).
\end{cuadrosolucion}

\textbf{Ejemplo R2.} Se quiere construir una caja abierta (sin tapa) a partir de una hoja cuadrada de $30$ cm de lado, recortando cuadrados iguales en las esquinas y doblando. ¿De qué tamaño deben ser los cuadrados recortados para maximizar el volumen?

\begin{cuadrosolucion}
  \textbf{Solución:} Sea $x$ el lado del cuadrado recortado ($0 < x < 15$). La base de la caja mide $30 - 2x$, y la altura es $x$.
  \begin{itemize}
    \item Volumen: $V(x) = x(30 - 2x)^2$.
    \item Derivada: $V'(x) = (30 - 2x)^2 + x \cdot 2(30 - 2x)(-2) = (30 - 2x)^2 - 4x(30 - 2x)$.
    \item Factorizando: $V'(x) = (30 - 2x)[(30 - 2x) - 4x] = (30 - 2x)(30 - 6x) = 6(15 - x)(5 - x)$.
    \item Ceros: $x = 15$ (fuera del dominio) y $x = 5$.
    \item $V''(5) = ...$ (verificar): $V''(x) = -24x + 240$? En $x = 5$: $V''(5) = -24(5) + 240 = 120 > 0$, espera, eso indicaría mínimo. Revisemos la derivada:
  \end{itemize}

  Rehagamos correctamente: $V(x) = x(30 - 2x)^2 = x(900 - 120x + 4x^2) = 900x - 120x^2 + 4x^3$.
  \begin{itemize}
    \item $V'(x) = 900 - 240x + 12x^2 = 12(x^2 - 20x + 75) = 12(x - 5)(x - 15)$.
    \item En $0 < x < 15$: $x = 5$ es el único punto crítico.
    \item $V''(x) = -240 + 24x$; $V''(5) = -240 + 120 = -120 < 0$: máximo local.
    \item $V(5) = 5(30 - 10)^2 = 5(20)^2 = 2000$ cm$^3$.
  \end{itemize}
  Se deben recortar cuadrados de $5$ cm para volumen máximo de $2000$ cm$^3$.
\end{cuadrosolucion}

\textbf{Ejemplo R3.} Encuentra las dimensiones del rectángulo de mayor área que puede inscribirse en una semicircunferencia de radio $5$, con la base sobre el diámetro.

\begin{cuadrosolucion}
  \textbf{Solución:} Colocamos el centro en el origen. El rectángulo tiene base $2x$ y altura $y = \sqrt{25 - x^2}$ (por la ecuación de la circunferencia $x^2 + y^2 = 25$). Dominio: $0 \le x \le 5$.
  \begin{itemize}
    \item Área: $A(x) = 2x\sqrt{25 - x^2}$.
    \item Derivada (producto): $A'(x) = 2\sqrt{25 - x^2} + 2x \cdot \frac{-x}{\sqrt{25 - x^2}} = \frac{2(25 - x^2) - 2x^2}{\sqrt{25 - x^2}} = \frac{50 - 4x^2}{\sqrt{25 - x^2}}$.
    \item Cero: $50 - 4x^2 = 0 \Rightarrow x^2 = 12{,}5 \Rightarrow x = \sqrt{12{,}5} = \frac{5\sqrt{2}}{2}$.
    \item Altura: $y = \sqrt{25 - 12{,}5} = \sqrt{12{,}5} = \frac{5\sqrt{2}}{2}$.
  \end{itemize}
  Las dimensiones son base $2x = 5\sqrt{2}$ y altura $y = \frac{5\sqrt{2}}{2}$.
\end{cuadrosolucion}

\newpage

\textbf{Ejemplo R4.} El costo de producir $x$ artículos es $C(x) = 0{,}1x^2 + 20x + 500$ dólares, y se venden a $50$ dólares cada uno. ¿Cuántos artículos maximizan la ganancia?

\begin{cuadrosolucion}
  \textbf{Solución:} Ingresos: $R(x) = 50x$. Ganancia: $P(x) = R(x) - C(x) = 50x - (0{,}1x^2 + 20x + 500) = -0{,}1x^2 + 30x - 500$.
  \begin{itemize}
    \item $P'(x) = -0{,}2x + 30 = 0 \Rightarrow x = 150$.
    \item $P''(x) = -0{,}2 < 0$: máximo.
  \end{itemize}
  Se deben producir y vender 150 artículos.
\end{cuadrosolucion}

\newpage

\subsection{Ejercicios propuestos}

\begin{enumerate}[leftmargin=*, resume]
  \item \facil{} Encuentra dos números positivos cuya suma sea $20$ y su producto sea máximo.

  \item \facil{} Encuentra dos números positivos cuyo producto sea $100$ y su suma sea mínima.

  \item \medio{} Un rectángulo tiene perímetro fijo de $40$ m. Encuentra sus dimensiones para que el área sea máxima.

  \item \medio{} Se lanza un proyectil hacia arriba y su altura es $h(t) = 80t - 16t^2$ pies. ¿Cuál es la altura máxima?

  \item \medio{} Un fabricante puede vender $100$ artículos a $10$ dólares cada uno. Por cada reducción de $0{,}50$ dólares en el precio, se venden 10 artículos más. ¿Cuál es el precio óptimo para maximizar el ingreso?

  \item \medio{} Se quiere construir un corral rectangular usando $120$ m de cerca, aprovechando un muro existente como uno de los lados. Encuentra las dimensiones que maximizan el área.

  \item \dificil{} Una caja cerrada con base cuadrada debe tener volumen $32$ cm$^3$. Encuentra las dimensiones que minimizan el área de superficie total.

  \item \dificil{} Encuentra el punto sobre la parábola $y = x^2$ más cercano al punto $(0, 3)$.

  \item \dificil{} Una escalera de $10$ m está apoyada contra una pared. Si la base se desliza alejándose de la pared a $2$ m/s, ¿qué tan rápido baja el extremo superior cuando la base está a $6$ m de la pared? (Usa el Teorema de Pitágoras.)

  \item \dificil{} Se vierte arena a razón de $10$ m$^3$/min formando un cono cuya altura es siempre igual al doble del radio. ¿A qué ritmo crece el radio cuando el radio es $5$ m? (Volumen del cono: $V = \frac{1}{3}\pi r^2 h$.)

  \item \dificil{} Un poste de $6$ m proyecta una sombra mientras el sol se pone. Un hombre de $1{,}80$ m camina alejándose del poste a $1{,}5$ m/s. ¿Qué tan rápido crece su sombra?

  \item \dificil{} Se desea hacer una caja abierta con base cuadrada y volumen $32$ cm$^3$ usando el menor material posible. ¿Cuáles son las dimensiones óptimas?
\end{enumerate}

\newpage

% ============================================================
\section{Teorema de Rolle y Teorema del Valor Medio}
% ============================================================

\subsection{Enunciados}

\begin{cuadroteoria}
  \textbf{Teorema de Rolle:} Si $f$ es continua en $[a, b]$, derivable en $(a, b)$, y $f(a) = f(b)$, entonces existe al menos un $c \in (a, b)$ tal que $f'(c) = 0$.

  \textbf{Teorema del Valor Medio (TVM):} Si $f$ es continua en $[a, b]$ y derivable en $(a, b)$, entonces existe al menos un $c \in (a, b)$ tal que:
  $$f'(c) = \frac{f(b) - f(a)}{b - a}$$
\end{cuadroteoria}

\begin{cuadrosolucion}
  \textbf{Interpretación geométrica del TVM:} existe un punto donde la tangente es paralela a la recta secante que une los extremos del intervalo.
\end{cuadrosolucion}

\subsection{Ejercicios resueltos}

\textbf{Ejemplo R1.} Verifica el Teorema de Rolle para $f(x) = x^2 - 4x$ en $[0, 4]$ y encuentra $c$.

\begin{cuadrosolucion}
  \textbf{Solución:}
  \begin{itemize}
    \item $f$ es polinomial: continua y derivable en todo $\mathbb{R}$.
    \item $f(0) = 0$; $f(4) = 16 - 16 = 0$. Se cumple $f(0) = f(4)$.
    \item $f'(x) = 2x - 4$. Igualamos a cero: $2x - 4 = 0 \Rightarrow x = 2$.
  \end{itemize}
  $c = 2$ pertenece a $(0, 4)$ y cumple $f'(c) = 0$.
\end{cuadrosolucion}

\textbf{Ejemplo R2.} Encuentra $c$ según el TVM para $f(x) = x^3$ en $[1, 2]$.

\begin{cuadrosolucion}
  \textbf{Solución:}
  \begin{itemize}
    \item Pendiente secante: $\dfrac{f(2) - f(1)}{2 - 1} = \dfrac{8 - 1}{1} = 7$.
    \item Derivada: $f'(x) = 3x^2$. Igualamos: $3c^2 = 7 \Rightarrow c^2 = \frac{7}{3} \Rightarrow c = \sqrt{\frac{7}{3}} \approx 1{,}53 \in (1, 2)$.
  \end{itemize}
\end{cuadrosolucion}

\textbf{Ejemplo R3.} Un automóvil acelera de $0$ a $100$ km/h en $10$ segundos. ¿Qué te dice el TVM?

\begin{cuadrosolucion}
  \textbf{Solución:} La velocidad promedio es $\dfrac{100 - 0}{10}$ km/h/s $= 10$ km/h/s. El TVM garantiza que en algún instante la velocidad instantánea (derivada de la posición) fue exactamente esa velocidad promedio. (Y de hecho, si la función de velocidad es continua, también se puede aplicar a la velocidad.)
\end{cuadrosolucion}

\newpage

\subsection{Ejercicios propuestos}

\begin{enumerate}[leftmargin=*, resume]
  \item \facil{} Verifica el Teorema de Rolle para $f(x) = x^2 - 1$ en $[-1, 1]$ y encuentra $c$.

  \item \facil{} Verifica el Teorema de Rolle para $f(x) = \sin x$ en $[0, \pi]$ y encuentra $c$.

  \item \facil{} Encuentra $c$ según el TVM para $f(x) = 2x + 1$ en $[0, 3]$.

  \item \facil{} Encuentra $c$ según el TVM para $f(x) = x^2$ en $[2, 4]$.

  \item \medio{} Verifica si se puede aplicar Rolle a $f(x) = x^2 - 4x + 3$ en $[1, 3]$ y halla $c$.

  \item \medio{} Encuentra $c$ según el TVM para $f(x) = \sqrt{x}$ en $[1, 4]$.

  \item \medio{} Un objeto se mueve con posición $s(t) = t^2 + 3t$ metros entre $t = 0$ y $t = 2$ s. ¿Cuál es la velocidad promedio? ¿En qué instante $t$ la velocidad instantánea coincide con la promedio?

  \item \medio{} Demuestra usando el TVM que si $f'(x) = 0$ para todo $x$, entonces $f$ es constante en el intervalo.

  \item \dificil{} Aplica el TVM a $f(x) = \ln x$ en $[1, e]$ y encuentra $c$.

  \item \dificil{} Sea $f$ continua y derivable con $f(1) = 2$ y $f(5) = 10$. ¿Qué puedes afirmar sobre $f'(c)$ para algún $c \in (1, 5)$?

  \item \dificil{} Prueba que la ecuación $x^3 - 3x + 1 = 0$ tiene al menos una raíz en $(0, 1)$ (pista: usa TVM sobre alguna función conveniente, o el TVI).

  \item \dificil{} Explica por qué el TVM no se puede aplicar a $f(x) = |x|$ en $[-1, 1]$.
\end{enumerate}

\newpage

% ============================================================
\section{Razones de Cambio Relacionadas}
% ============================================================

\subsection{Metodología}

\begin{cuadroteoria}
  \textbf{Metodología de resolución:}
  \begin{enumerate}[leftmargin=*]
    \item \textbf{Dibuja un esquema} y asigna variables a las cantidades que varían con el tiempo $t$.
    \item \textbf{Identifica las razones conocidas y desconocidas} (derivadas respecto a $t$, como $\frac{dx}{dt}$, $\frac{dV}{dt}$).
    \item \textbf{Escribe una ecuación} que relacione las variables (Pitágoras, fórmulas de volumen/área, triángulos semejantes).
    \item \textbf{Deriva implícitamente} ambos lados de la ecuación respecto al tiempo $t$ (regla de la cadena).
    \item \textbf{Sustituye} los valores conocidos en el instante dado y despeja la razón buscada.
  \end{enumerate}
\end{cuadroteoria}

\begin{cuadroadvertencia}
  \textbf{¡Cuidado!} NUNCA sustituyas valores numéricos de variables cambiantes ANTES de derivar. Sustituye únicamente DESPUÉS de haber derivado respecto a $t$.
\end{cuadroadvertencia}

\subsection{Ejercicios resueltos}

\textbf{Ejemplo R1.} Una escalera de $10$ m de largo está apoyada contra una pared vertical. Si la base de la escalera se desliza alejándose de la pared a razón de $2$ m/s, ¿con qué rapidez baja el extremo superior cuando la base está a $6$ m de la pared?

\begin{cuadrosolucion}
  \textbf{Solución:}
  \begin{enumerate}
    \item Variables: $x(t) =$ distancia de la base a la pared, $y(t) =$ altura en la pared.
    \item Datos: $\dfrac{dx}{dt} = 2$ m/s. Queremos $\dfrac{dy}{dt}$ cuando $x = 6$.
    \item Relación (Pitágoras): $x^2 + y^2 = 10^2 = 100$.
    \item Derivamos respecto a $t$: $2x \dfrac{dx}{dt} + 2y \dfrac{dy}{dt} = 0 \Rightarrow x \dfrac{dx}{dt} + y \dfrac{dy}{dt} = 0$.
    \item Para $x = 6$: $y = \sqrt{100 - 36} = 8$ m.
    \item Sustituimos: $6(2) + 8 \dfrac{dy}{dt} = 0 \Rightarrow 12 + 8 \dfrac{dy}{dt} = 0 \Rightarrow \dfrac{dy}{dt} = -\dfrac{12}{8} = -1{,}5$ m/s.
  \end{enumerate}
  El extremo superior desciende a $1{,}5$ m/s (el signo negativo indica que $y$ disminuye).
\end{cuadrosolucion}

\textbf{Ejemplo R2.} Se vierte agua en un tanque cónico recetáculo (con el vértice hacia abajo) de radio $4$ m y altura $10$ m a razón de $2$ m$^3$/min. ¿Con qué rapidez sube el nivel del agua cuando la profundidad es de $5$ m?

\begin{cuadrosolucion}
  \textbf{Solución:}
  \begin{enumerate}
    \item Variables: $r =$ radio de la superficie del agua, $h =$ profundidad del agua, $V =$ volumen de agua.
    \item Datos: $\dfrac{dV}{dt} = 2$ m$^3$/min. Queremos $\dfrac{dh}{dt}$ cuando $h = 5$.
    \item Volumen del cono: $V = \dfrac{1}{3}\pi r^2 h$.
    \item Triángulos semejantes: $\dfrac{r}{h} = \dfrac{4}{10} \Rightarrow r = \dfrac{2}{5}h$.
    \item Expresamos $V$ solo en función de $h$: $V = \dfrac{1}{3}\pi \left(\dfrac{2}{5}h\right)^2 h = \dfrac{4\pi}{75} h^3$.
    \item Derivamos respecto a $t$: $\dfrac{dV}{dt} = \dfrac{4\pi}{75} \cdot 3h^2 \dfrac{dh}{dt} = \dfrac{4\pi}{25} h^2 \dfrac{dh}{dt}$.
    \item Sustituimos $h = 5$ y $\dfrac{dV}{dt} = 2$:
          $$2 = \dfrac{4\pi}{25}(25) \dfrac{dh}{dt} = 4\pi \dfrac{dh}{dt} \Rightarrow \dfrac{dh}{dt} = \dfrac{2}{4\pi} = \dfrac{1}{2\pi} \approx 0{,}159 \text{ m/min}.$$
  \end{enumerate}
\end{cuadrosolucion}

\newpage

\subsection{Ejercicios propuestos}

\begin{enumerate}[leftmargin=*, resume]
  \item \facil{} Un globo esférico se infla a razón de $100$ cm$^3$/s. ¿Con qué rapidez aumenta el radio cuando el radio mide $5$ cm? ($V = \frac{4}{3}\pi r^3$).
  \item \facil{} Las aristas de un cubo aumentan a razón de $3$ cm/s. ¿Con qué rapidez aumenta el volumen del cubo cuando cada arista mide $10$ cm?
  \item \medio{} Un bote es arrastrado hacia un muelle por medio de un cable atado a la proa y pasado por una polea situada en el muelle a $3$ m por encima de la proa. Si el cable se recoge a $1$ m/s, ¿con qué rapidez se aproxima el bote al muelle cuando la distancia del cable es $5$ m?
  \item \medio{} Un hombre de $1{,}80$ m de estatura camina alejándose de un farol de $6$ m de altura a razón de $1{,}5$ m/s. ¿Con qué rapidez crece su sombra?
  \item \dificil{} Dos automóviles parten del mismo punto: uno viaja hacia el norte a $60$ km/h y el otro hacia el este a $80$ km/h. ¿Con qué rapidez aumenta la distancia entre ellos $2$ horas después de la salida?
  \item \dificil{} Se vierte arena sobre el suelo a razón de $10$ m$^3$/min formando un montón cónico cuya altura es siempre igual al diámetro de la base. ¿Con qué rapidez aumenta la altura cuando el montón mide $4$ m de alto?
\end{enumerate}

\newpage

% ============================================================
\section{Análisis Completo de una Gráfica}
% ============================================================

\subsection{Concavidad y puntos de inflexión}

\begin{cuadroteoria}
  \textbf{Concavidad:}
  \begin{itemize}[leftmargin=*]
    \item Si $f''(x) > 0$ en un intervalo: cóncava hacia arriba (la curva se abre como una taza).
    \item Si $f''(x) < 0$ en un intervalo: cóncava hacia abajo (la curva se abre como una montaña).
  \end{itemize}

  \textbf{Punto de inflexión:} punto donde la concavidad cambia (donde $f''(x) = 0$ o no existe, y cambia de signo).
\end{cuadroteoria}

\begin{cuadropasos}
  \textbf{Pasos para graficar una función con análisis:}
  \begin{enumerate}
    \item Determina el \textbf{dominio}.
    \item Encuentra las \textbf{intersecciones} con los ejes.
    \item Analiza \textbf{simetrías} (par/impar).
    \item Encuentra las \textbf{asíntotas} (verticales, horizontales, oblicuas).
    \item Calcula $f'$ y determina \textbf{intervalos de crecimiento/decrecimiento} y extremos.
    \item Calcula $f''$ y determina \textbf{concavidad} y puntos de inflexión.
    \item Esboza la gráfica con toda la información.
  \end{enumerate}
\end{cuadropasos}

\subsection{Ejercicios resueltos}

\textbf{Ejemplo R1.} Determina intervalos de crecimiento, concavidad y puntos de inflexión de $f(x) = x^3 - 6x^2 + 9x$.

\begin{cuadrosolucion}
  \textbf{Solución:}
  \begin{enumerate}
    \item Dominio: $\mathbb{R}$.
    \item $f'(x) = 3x^2 - 12x + 9 = 3(x^2 - 4x + 3) = 3(x - 1)(x - 3)$.
          \begin{itemize}
            \item $f'(x) > 0$ en $(-\infty, 1) \cup (3, \infty)$: creciente.
            \item $f'(x) < 0$ en $(1, 3)$: decreciente.
            \item $x = 1$: máximo local ($f(1) = 4$). $x = 3$: mínimo local ($f(3) = 0$).
          \end{itemize}
    \item $f''(x) = 6x - 12 = 6(x - 2)$.
          \begin{itemize}
            \item $f''(x) < 0$ en $(-\infty, 2)$: cóncava hacia abajo.
            \item $f''(x) > 0$ en $(2, \infty)$: cóncava hacia arriba.
            \item Punto de inflexión en $x = 2$: $f(2) = 8 - 24 + 18 = 2$. Punto $(2, 2)$.
          \end{itemize}
  \end{enumerate}
\end{cuadrosolucion}

\textbf{Ejemplo R2.} Analiza y grafica $f(x) = \dfrac{x^2}{x - 1}$.

\begin{cuadrosolucion}
  \textbf{Solución:}
  \begin{enumerate}
    \item Dominio: $x \neq 1$.
    \item Interceptos: $x = 0$, $y = 0$. Punto $(0, 0)$.
    \item Asíntotas: A.V. en $x = 1$ (el denominador se anula y el numerador no). A.O.: dividiendo $x^2$ entre $x - 1$: $x + 1 + \frac{1}{x-1}$, así que $y = x + 1$ es asíntota oblicua.
    \item $f'(x) = \dfrac{2x(x - 1) - x^2}{(x - 1)^2} = \dfrac{x^2 - 2x}{(x - 1)^2} = \dfrac{x(x - 2)}{(x - 1)^2}$.
          \begin{itemize}
            \item Signo: positiva en $(-\infty, 0) \cup (2, \infty)$; negativa en $(0, 1) \cup (1, 2)$.
            \item $x = 0$: máximo local $f(0) = 0$. $x = 2$: mínimo local $f(2) = 4$.
          \end{itemize}
    \item $f''(x) = \dfrac{2}{(x - 1)^3}$ (verificar): es negativa en $(-\infty, 1)$ y positiva en $(1, \infty)$. Concavidad hacia abajo a la izquierda de 1, hacia arriba a la derecha.
  \end{enumerate}
\end{cuadrosolucion}

\newpage

\subsection{Ejercicios propuestos}

\begin{enumerate}[leftmargin=*, resume]
  \item \facil{} Determina la concavidad de $f(x) = x^2$ en todo su dominio.

  \item \facil{} Determina la concavidad de $f(x) = -x^2 + 3x$ en todo su dominio.

  \item \facil{} Encuentra los puntos de inflexión de $f(x) = x^3$.

  \item \facil{} Determina los intervalos de crecimiento de $f(x) = x^2 - 4x$.

  \item \medio{} Para $f(x) = x^3 - 3x$:
        \begin{enumerate}[label=\alph*)]
          \item Encuentra los puntos críticos y clasifícalos.
          \item Intervalos de crecimiento/decrecimiento.
          \item Intervalos de concavidad y puntos de inflexión.
        \end{enumerate}

  \item \medio{} Para $f(x) = \dfrac{1}{x}$:
        \begin{enumerate}[label=\alph*)]
          \item Determina los intervalos de crecimiento/decrecimiento.
          \item Determina la concavidad en cada intervalo del dominio.
          \item Identifica las asíntotas.
        \end{enumerate}

  \item \medio{} Encuentra los puntos de inflexión de $f(x) = x^4 - 6x^2$.

  \item \medio{} Analiza la gráfica de $f(x) = x e^{-x}$: dominio, extremos, crecimiento, concavidad y asíntota horizontal.

  \item \dificil{} Analiza y grafica $f(x) = \dfrac{2x}{x^2 - 1}$: dominio, interceptos, asíntotas, extremos, concavidad.

  \item \dificil{} Determina el valor de $a$ para que $f(x) = x^3 - ax^2 + 3$ tenga un punto de inflexión en $x = 1$.

  \item \dificil{} Analiza $f(x) = \ln(x^2 + 1)$: crecimiento, concavidad, extremos y puntos de inflexión.

  \item \dificil{} Una función $f$ es derivable dos veces. Si $f''(x) = 6x - 4$, $f'(1) = 2$ y $f(1) = 3$, determina $f(x)$.
\end{enumerate}

\newpage

% ============================================================
\section{Clave de Respuestas}
% ============================================================

\subsection*{Sección 1: Máximos y Mínimos (Ejercicios 1--12)}

\begin{enumerate}[leftmargin=*, label=\textbf{\arabic*.}]
  \item $f'(x) = 2x - 6$, punto crítico: $x = 3$.
  \item $f'(x) = 4x^3 - 8x = 4x(x^2 - 2)$. Puntos críticos: $x = 0$, $x = \pm\sqrt{2}$.
  \item Mínimo absoluto en $(0, 0)$: $0$. Máximo absoluto en $(-2, 4)$: $4$.
  \item Mínimo absoluto en $(0, 1)$: $1$. Máximo absoluto en $(4, 5)$: $5$.
  \item $f'(x) = 6x^2 - 6x - 36 = 6(x^2 - x - 6) = 6(x - 3)(x + 2)$. Ceros: $x = 3$ y $x = -2$. Tabla de signos: $f'$ positiva en $(-\infty, -2)$, negativa en $(-2, 3)$, positiva en $(3, \infty)$. En $x = -2$: máximo local. En $x = 3$: mínimo local.
  \item $f'(x) = 3x^2 - 3$, cero en $x = \pm 1$. En $[0, 3]$ solo $x = 1$ está adentro. $f(0) = 1$, $f(1) = -1$, $f(3) = 19$. Mínimo: $-1$ en $x = 1$; máximo: $19$ en $x = 3$.
  \item $f'(x) = \dfrac{1 - x^2}{(x^2 + 1)^2}$. Ceros: $x = \pm 1$.
  \item $f'(x) = \cos x$. Ceros en $x = \frac{\pi}{2}$ (máximo, $f = 1$) y $x = \frac{3\pi}{2}$ (mínimo, $f = -1$). Además, los extremos $0$ y $2\pi$ dan $f = 0$.
  \item $f'(x) = 3ax^2 - 12x$. En $x = 2$: $f'(2) = 3a(4) - 24 = 12a - 24 = 0 \Rightarrow a = 2$.
  \item $f'(x) = \sqrt{9 - x^2} + x \cdot \frac{-x}{\sqrt{9-x^2}} = \frac{9 - 2x^2}{\sqrt{9-x^2}}$. Ceros: $x = \pm \frac{3\sqrt{2}}{2}$. Evaluando en esos y en los extremos: máximo $\frac{9}{2}$ en $x = \pm \frac{3\sqrt{2}}{2}$; mínimo $0$ en los extremos $x = \pm 3$.
  \item $f'(x) = \dfrac{(x + 1) - (x - 1)}{(x + 1)^2} = \dfrac{2}{(x + 1)^2} > 0$ siempre, creciente. Mínimo en $x = 0$: $f(0) = -1$. Máximo en $x = 4$: $f(4) = \frac{3}{5}$.
  \item $f'(x) = 3x^2 + 1 > 0$ siempre, no hay puntos críticos (la derivada nunca se anula). $f$ es estrictamente creciente en todo $\mathbb{R}$.
\end{enumerate}

\subsection*{Sección 2: Optimización (Ejercicios 13--24)}

\begin{enumerate}[leftmargin=*, label=\textbf{\arabic*.}]
  \setcounter{enumi}{12}
  \item Ambos números $= 10$ (producto máximo $100$).
  \item Ambos números $= 10$ (suma mínima $20$).
  \item Cuadrado de lado $10$ m (área máxima $100$ m$^2$).
  \item $h'(t) = 80 - 32t = 0 \Rightarrow t = 2{,}5$ s. $h(2{,}5) = 80(2{,}5) - 16(6{,}25) = 200 - 100 = 100$ pies.
  \item Si $x$ = reducción en dólares (múltiplos de $0{,}5$): precio $= 10 - x$, cantidad $= 100 + 20x$. Ingreso $R = (10 - x)(100 + 20x)$. Derivando: $R' = -100 + 200 - 40x = 100 - 40x = 0 \Rightarrow x = 2{,}5$. Precio óptimo: $7{,}50$ dólares.
  \item Con muro como lado, largo $= 120 - 2x$ (donde $x$ es el ancho). Área $= x(120 - 2x)$. $A' = 120 - 4x = 0 \Rightarrow x = 30$. Dimensiones: $30 \times 60$ m.
  \item Base cuadrada $x$, altura $h = \frac{32}{x^2}$. Área total: $2x^2 + 4xh = 2x^2 + \frac{128}{x}$. $A' = 4x - \frac{128}{x^2} = 0 \Rightarrow 4x^3 = 128 \Rightarrow x = \sqrt[3]{32} = 2\sqrt[3]{4}$. Altura $h = \frac{32}{x^2} = \frac{32}{(2\sqrt[3]{4})^2}$. Simplificado: $x = 2\sqrt[3]{4}$, $h = \frac{32}{4 \sqrt[3]{16}} = \frac{8}{\sqrt[3]{16}} = \frac{8}{2\sqrt[3]{2}} = \frac{4}{\sqrt[3]{2}} = 2\sqrt[3]{4}$. ¡Es un cubo! (base $= h$).
  \item Punto $(a, a^2)$. Distancia al cuadrado al cuadrado al cuadrado al cuadrd: $D = a^2 + (a^2 - 3)^2$. $D' = 2a + 2(a^2 - 3)(2a) = 2a + 4a(a^2 - 3) = 2a(1 + 2a^2 - 6) = 2a(2a^2 - 5) = 0$. $a = 0$ o $a = \pm\sqrt{2{,}5}$. Los puntos más cercanos: $\left(\sqrt{2{,}5},\ 2{,}5\right)$, $\left(-\sqrt{2{,}5},\ 2{,}5\right)$. (El punto $(0,0)$ es un máximo local de la distancia).
  \item Sea $x$ = distancia de la base a la pared, $y$ = altura. $x^2 + y^2 = 100$. Derivando: $2x\frac{dx}{dt} + 2y\frac{dy}{dt} = 0$. En $x = 6$, $y = 8$, $\frac{dx}{dt} = 2$: $2(6)(2) + 2(8)\frac{dy}{dt} = 0 \Rightarrow \frac{dy}{dt} = -\frac{24}{16} = -1{,}5$ m/s (la altura baja a 1,5 m/s).
  \item $V = \frac{1}{3}\pi r^2(2r) = \frac{2}{3}\pi r^3$. Derivando: $\frac{dV}{dt} = 2\pi r^2 \frac{dr}{dt}$. Con $\frac{dV}{dt} = 10$ y $r = 5$: $10 = 2\pi(25)\frac{dr}{dt} \Rightarrow \frac{dr}{dt} = \frac{10}{50\pi} = \frac{1}{5\pi}$ m/min $\approx 0{,}064$ m/min.
  \item Por triángulos semejantes: $\frac{6}{x + s} = \frac{1{,}8}{s}$. Entonces $6s = 1{,}8x + 1{,}8s \Rightarrow 4{,}2s = 1{,}8x \Rightarrow s = \frac{3}{7}x$. Derivando: $\frac{ds}{dt} = \frac{3}{7}\frac{dx}{dt} = \frac{3}{7}(1{,}5) = \frac{4{,}5}{7} \approx 0{,}643$ m/s.
  \item Similar al anterior pero con base cuadrada $x$ y volumen fijo $32$: $h = \frac{32}{x^2}$. Área total (sin tapa) $= x^2 + 4xh = x^2 + \frac{128}{x}$. $A' = 2x - \frac{128}{x^2} = 0 \Rightarrow 2x^3 = 128 \Rightarrow x = 4$ cm, $h = \frac{32}{16} = 2$ cm.
\end{enumerate}

\subsection*{Sección 3: Rolle y TVM (Ejercicios 25--36)}

\begin{enumerate}[leftmargin=*, label=\textbf{\arabic*.}]
  \setcounter{enumi}{24}
  \item $f(-1) = 0$, $f(1) = 0$. $f'(x) = 2x = 0 \Rightarrow c = 0$.
  \item $f(0) = 0$, $f(\pi) = 0$. $f'(x) = \cos x = 0 \Rightarrow c = \frac{\pi}{2}$.
  \item $f$ es lineal con pendiente 2, así que $c$ puede ser cualquier punto en $(0, 3)$ (por ejemplo $c = 1$).
  \item $\frac{f(4) - f(2)}{4 - 2} = \frac{16 - 4}{2} = 6$. $f'(x) = 2x = 6 \Rightarrow c = 3$.
  \item $f(1) = 1 - 4 + 3 = 0$, $f(3) = 9 - 12 + 3 = 0$. $f$ es continua y derivable. $f'(x) = 2x - 4 = 0 \Rightarrow c = 2 \in (1, 3)$.
  \item Pendiente secante: $\frac{2 - 1}{4 - 1} = \frac{1}{3}$. $f'(x) = \frac{1}{2\sqrt{x}} = \frac{1}{3} \Rightarrow 2\sqrt{c} = 3 \Rightarrow c = \frac{9}{4}$.
  \item Posición promedio: $\frac{s(2) - s(0)}{2} = \frac{10 - 0}{2} = 5$ m/s. $v(t) = s'(t) = 2t + 3 = 5 \Rightarrow t = 1$ s.
  \item Por el TVM, para cualquier $a < b$: $f'(c) = \frac{f(b) - f(a)}{b - a} = 0 \Rightarrow f(b) = f(a)$. Como $a$ y $b$ son arbitrarios, $f$ es constante.
  \item $\frac{f(e) - f(1)}{e - 1} = \frac{1 - 0}{e - 1} = \frac{1}{e - 1}$. $f'(x) = \frac{1}{x} = \frac{1}{e - 1} \Rightarrow c = e - 1 \approx 1{,}718 \in (1, e)$.
  \item TVM: existe $c \in (1, 5)$ con $f'(c) = \frac{10 - 2}{5 - 1} = 2$.
  \item Sea $f(x) = x^3 - 3x + 1$. $f(0) = 1$, $f(1) = -1$. Por TVI existe al menos una raíz en $(0,1)$. (También se puede argumentar con el TVM sobre una primitiva.)
  \item $f(x) = |x|$ no es derivable en $x = 0 \in (-1, 1)$, que es un punto interior del intervalo. No se cumplen las hipótesis del TVM.
\end{enumerate}

\subsection*{Sección 4: Razones de Cambio Relacionadas (Ejercicios 37--42)}

\begin{enumerate}[leftmargin=*, label=\textbf{\arabic*.}]
  \setcounter{enumi}{36}
  \item $V = \frac{4}{3}\pi r^3 \Rightarrow \frac{dV}{dt} = 4\pi r^2 \frac{dr}{dt} \Rightarrow 100 = 4\pi(25)\frac{dr}{dt} \Rightarrow \frac{dr}{dt} = \frac{1}{\pi} \approx 0{,}318$ cm/s.
  \item $V = x^3 \Rightarrow \frac{dV}{dt} = 3x^2 \frac{dx}{dt} = 3(100)(3) = 900$ cm$^3$/s.
  \item Sea $x$ la distancia al muelle y $z$ la longitud del cable: $x^2 + 3^2 = z^2$. Derivando: $2x\frac{dx}{dt} = 2z\frac{dz}{dt}$. Con $z = 5 \Rightarrow x = 4$. $4\frac{dx}{dt} = 5(-1) \Rightarrow \frac{dx}{dt} = -1{,}25$ m/s. Se aproxima a $1{,}25$ m/s.
  \item Triángulos semejantes: $\frac{6}{x+s} = \frac{1{,}8}{s} \Rightarrow s = \frac{3}{7}x \Rightarrow \frac{ds}{dt} = \frac{3}{7}\frac{dx}{dt} = \frac{3}{7}(1{,}5) = \frac{4{,}5}{7} \approx 0{,}643$ m/s.
  \item Distancia $D = \sqrt{x^2 + y^2}$. A las 2 h: $x = 160$ km, $y = 120$ km $\Rightarrow D = 200$ km. Derivando: $\frac{dD}{dt} = \frac{x\frac{dx}{dt} + y\frac{dy}{dt}}{D} = \frac{160(80) + 120(60)}{200} = \frac{12800 + 7200}{200} = \frac{20000}{200} = 100$ km/h.
  \item Diámetro $= 2r = h \Rightarrow r = \frac{h}{2}$. $V = \frac{1}{3}\pi r^2 h = \frac{\pi}{12}h^3$. Derivando: $\frac{dV}{dt} = \frac{\pi}{4}h^2 \frac{dh}{dt} \Rightarrow 10 = \frac{\pi}{4}(16)\frac{dh}{dt} = 4\pi \frac{dh}{dt} \Rightarrow \frac{dh}{dt} = \frac{10}{4\pi} = \frac{5}{2\pi} \approx 0{,}796$ m/min.
\end{enumerate}

\subsection*{Sección 5: Análisis de Gráficas (Ejercicios 43--54)}

\begin{enumerate}[leftmargin=*, label=\textbf{\arabic*.}]
  \setcounter{enumi}{42}
  \item $f''(x) = 2 > 0$: cóncava hacia arriba en todo $\mathbb{R}$.
  \item $f''(x) = -2 < 0$: cóncava hacia abajo en todo $\mathbb{R}$.
  \item $f''(x) = 6x = 0 \Rightarrow x = 0$; cambia de signo (negativa a la izquierda, positiva a la derecha). Punto de inflexión en $(0, 0)$.
  \item $f'(x) = 2x - 4$; creciente en $(2, \infty)$, decreciente en $(-\infty, 2)$.
  \item a) $f'(x) = 3x^2 - 3$; puntos críticos $x = \pm 1$. $x = -1$ máximo, $x = 1$ mínimo. b) Creciente en $(-\infty, -1) \cup (1, \infty)$; decreciente en $(-1, 1)$. c) $f''(x) = 6x$; cóncava abajo en $(-\infty, 0)$, arriba en $(0, \infty)$; punto de inflexión $(0, 0)$.
  \item a) $f'(x) = -\frac{1}{x^2} < 0$: decreciente en $(-\infty, 0)$ y en $(0, \infty)$. b) $f''(x) = \frac{2}{x^3}$: cóncava abajo en $(-\infty, 0)$, cóncava arriba en $(0, \infty)$. c) Asíntotas: $x = 0$ (vertical) e $y = 0$ (horizontal).
  \item $f''(x) = 12x^2 - 12 = 12(x^2 - 1) = 12(x - 1)(x + 1)$. Puntos de inflexión: $x = \pm 1$ (ambos con cambio de signo).
  \item Dominio $\mathbb{R}$. $f'(x) = e^{-x}(1 - x)$: creciente en $(-\infty, 1)$, decreciente en $(1, \infty)$; máximo en $x = 1$, $f(1) = e^{-1}$. $f''(x) = e^{-x}(x - 2)$: cóncava abajo en $(-\infty, 2)$, arriba en $(2, \infty)$; punto de inflexión $(2, 2e^{-2})$. Asíntota horizontal $y = 0$ cuando $x \to \infty$.
  \item Dominio: $x \neq \pm 1$. Intercepto $(0, 0)$. Asíntotas verticales: $x = 1$ y $x = -1$. Asíntota horizontal: $y = 0$. $f'(x) = \frac{-2x^2 - 2}{(x^2 - 1)^2} < 0$ siempre: decreciente en cada intervalo. $f''$ analiza concavidad; la información completa se usa para el esbozo.
  \item $f''(x) = 6x - 2a$. En $x = 1$: $6 - 2a = 0 \Rightarrow a = 3$ (y además cambia de signo, así que es punto de inflexión).
  \item $f'(x) = \frac{2x}{x^2 + 1}$: decreciente en $(-\infty, 0)$, creciente en $(0, \infty)$; mínimo en $(0, 0)$. $f''(x) = \frac{-2x^2 + 2}{(x^2 + 1)^2} = \frac{2(1 - x^2)}{(x^2 + 1)^2}$: cóncava arriba en $(-1, 1)$, abajo en $(-\infty, -1) \cup (1, \infty)$; puntos de inflexión en $x = \pm 1$.
  \item $f(x) = \displaystyle\int (6x - 4)\,dx = 3x^2 - 4x + C$. Con $f'(1) = 2$: $3 - 4 + C = 2 \Rightarrow C = 3$. Entonces $f(x) = 3x^2 - 4x + 3$. Integrando: $\int(3x^2 - 4x + 3)\,dx = x^3 - 2x^2 + 3x + D$. Con $f(1) = 3$: $1 - 2 + 3 + D = 3 \Rightarrow D = 1$. Por tanto $f(x) = x^3 - 2x^2 + 3x + 1$.
\end{enumerate}

\end{document}